%% HAND-WRITTEN fragment -- NOT generated by maestro2tex, placed by it via the
%% `order` list in configs/anatop_tb_pixel_corners.json (block Pixel). The
%% engine has no signal-vs-signal (XY) figure kind, only signal-vs-native-
%% sweep-axis, so this true current-vs-potential voltammogram is built from a
%% merged .dat instead of a generated fig_*.tex.
%%
%% Data: data/px_cv_voltammogram_tt.dat, written by
%% skill/px_cv_merge.py from the ordinary time-domain CSVs
%% data/raw/px_cv_tran__tt__{cv_x,cv_i}.csv that configs/anatop_tb_pixel_cv.json
%% (block PixelCV) extracts from
%% simruns/tb_anatop_pixel_cv_1M/.../Ocean.0 (rf=35k, rct_we=1M, cdl_we=1u,
%% typical process point, 37 C). Re-run px_cv_merge.py after any re-extraction
%% of PixelCV. Palette colour mtxA matches the other Pixel current-vs-x traces.
\providecommand{\MaestroRoot}{include/analog/}
\begin{figure}[htbp]
  \centering
  {\pgfplotsset{compat=1.18}%
  \begin{tikzpicture}
    \begin{axis}[
      width=0.68\linewidth, height=6.0cm,
      xlabel={RE potential~[\si{\volt}]}, ylabel={WE current~[\si{\micro\ampere}]},
      grid=both,
      major grid style={line width=0.2pt, draw=black!14},
      minor grid style={line width=0.2pt, draw=black!7},
      minor tick num=1,
      axis line style={line width=0.4pt, draw=black!45},
      tick style={line width=0.4pt, draw=black!45},
      tick align=outside,
      label style={font=\small}, tick label style={font=\small},
      enlarge x limits=0.05, enlarge y limits=0.08,
      every axis plot/.append style={line join=round, no marks},
    ]
      \addplot[mtxA, line width=1.1pt, solid] table {\MaestroRoot data/px_cv_voltammogram_tt.dat};
    \end{axis}
  \end{tikzpicture}}
  \caption[{CV voltammogram of the pixel system.}]{Working-electrode current against commanded reference-electrode potential over one triangular voltammetry sweep, typical process point, \SI{37}{\celsius}, $R_f = \SI{35}{\kilo\ohm}$, $R_{\mathrm{ct}} = \SI{1}{\mega\ohm}$ electrode. The trace is an ellipse-like loop rather than a single-valued curve because the double-layer capacitance $C_{\mathrm{dl}}$ carries a current proportional to $\mathrm{d}v_{\mathrm{re}}/\mathrm{d}t$ in addition to the resistive charge-transfer current; the two straight sweep segments (RE potential linear in time, Figure~\ref{fig:px-cv-potential}) are why the loop closes into two nearly-parallel branches rather than a smooth curve. The current stays unclipped over the whole sweep (Table~\ref{tab:px-cv-stats}).}
  \label{fig:px-cv-voltammogram}
\end{figure}
